\chapter{Systèmes inverses, représentation de Weyl--Heisenberg et quotient prédictif}
\label{chap:numero-heisenberg-d108}
\index{nombre HMT}
\index{Heisenberg!groupe fini}

La construction distingue trois niveaux : le catalogue fini des germes,
l'espace des états qui conserve la retenue et l'orientation, et la limite
inverse des prolongements compatibles. Cette stratification permet de
formaliser la contraction des intervalles cylindriques jusqu'à un point réel,
la persistance des données de trajectoire dans une même fibre d'évaluation et
la représentation finie de Weyl--Heisenberg associée à la forme alternée
du tore nonadique qui porte les observables additif et multiplicatif.

Les symboles \(V_t\), \(X_t\) et \(\rho_{t+1,t}\) utilisés ici sont les
mêmes ensembles d'états locaux, préfixes admissibles et troncatures
introduits dans le \cref{chap:numero-medicion}, avec le changement purement
notationnel \(t=m\) et \(\rho_{n,m}=p_{n,m}\). Cette section explicite leurs
coordonnées et n'introduit pas de second système inverse.

\section{Projection observable et espace complet des états du TPK}

La projection observable du système de transport de phase (TPK) est
l'application finie de codage
\[
104\,976\longrightarrow468\longrightarrow243.
\]
qui classe les germes, signatures et mots \(w_6\). L'espace complet des
états du TPK conserve en outre
\[
\begin{gathered}
\text{orientation bidirectionnelle},\quad
\text{état de Hensel},\quad
\text{cylindre et retenue},\\
\text{signature conjointe de frontière},\quad
\text{phase nonadique visible }g_{\rm vis}(m).
\end{gathered}
\]
Deux candidats ayant une même signature locale peuvent appartenir à des classes distinctes
lorsque la frontière suivante est conservée. Nous n'introduisons pas de second tuple
d'état. Pour
\[
x_m=(q_m,o_m,h_m,c_m,\sigma_m,C_m,b_m,\lambda_m)\in\XTPK^{(m)}
\]
nous utilisons la carte de coordonnées
\[
\chi_m(x_m)=
\bigl(
B_m,h_m,(C_m,b_m),\sigma_m,
o_m,\mu(q_m),g_{\rm vis}(m),\lambda_m
\bigr),
\]
avec
\[
B_m=(u_m^\pi,u_m^e,u_m^\varphi)\in(\mathbb F_3^6)^3.
\]
Ici \(h_m\) est l'extension de Hensel ; \((C_m,b_m)\), le cylindre et sa
frontière ; \(o_m,\mu(q_m)\), l'orientation et l'observable bidirectionnel ; et
\[
g_{\rm vis}(m)=1+((m-1)\bmod9)\in\{1,\ldots,9\},
\]
l'étiquette visible de phase nonadique. La coordonnée de classe correspondante
est \(g_0(\theta_m)\in\ZZ/9\ZZ\), avec \(\theta_m=[m-1]_{108}\). La signature conjointe
est la coordonnée déjà typée par le noyau de phase,
\[
\sigma_m=\operatorname{Sig}_{q_m}(h_m,c_m,\lambda_m)
=(q_\partial,a_\partial,c_\partial,
\operatorname{colw},\rho_W,r_\partial).
\]
La signature conjointe ne se réduit pas à trois signatures de canal indépendantes, car la
sélection associée à l'involution \(\pi/e\) et à l'axe \(\varphi\) dépend de
données corrélées. Deux nœuds ayant le même \(B_m\) représentent des états
distincts lorsqu'ils diffèrent par l'extension, le cylindre, la retenue, la signature, l'orientation
ou le registre de transport.

\subsection{Compatibilité exacte des cylindres sous troncature et prolongement nonadique}

Soit \(N\) l'entier d'un préfixe ternaire de longueur \(L\), et soit \(D\)
l'entier obtenu en concaténant ses \(b\) triades décimales. Nous définissons
\[
A=N1000^b-D3^L,
\qquad
B=(N+1)1000^b-D3^L.
\]
Les cylindres se coupent si et seulement si
\[
\boxed{A<3^L,\qquad B>0.}
\]
En ajoutant un bloc ternaire \(u\in\{0,\ldots,728\}\), on a
\[
N'=729N+u.
\]
S'il n'apparaît pas de nouvelle triade décimale,
\[
A'=729A+u1000^b.
\]
Si \(\delta\) apparaît, de sorte que \(D'=1000D+\delta\), alors
\[
A'=729000A+u1000^{b+1}-\delta\,729\,3^L.
\]
Si \(\Delta b\in\{0,1\}\) enregistre l'apparition de cette nouvelle triade, dans
les deux cas
\[
B'=A'+1000^{b+\Delta b}.
\]
La transition, la paire admissible \((u,\delta)\) étant fixée, ne consulte pas de valeur
cible. La filtration conjointe par intersection de cylindres et retour
nonadique agit à un niveau différent : elle sélectionne les paires de frontière qui
admettent un prolongement compatible.

\subsection{Registre cyclique de 108 étapes et sa projection réduite}

Nous appelons \emph{registre cyclique de 108 étapes} la suite de
\(108\) pas regroupés en douze secteurs de neuf pas, avec la convention
d'indice fixée dans le \cref{chap:tpk-definicion}.

Soit \(\XTPK^{[108]}\subseteq\XTPK\) le domaine des états complets dont le
registre \(\lambda\) contient les \(108\) pas déclarés. Notons
\(\mathcal C_{108}^{\rm red}\) l'ensemble des registres qui résultent de la
conservation des seuls champs réduits, et \(\mathcal U_{12}^{\rm sgn}\)
l'ensemble des dodécaphases signées admissibles. L'instance distingue deux applications
déclarées :
\[
F:\XTPK^{[108]}\longrightarrow\mathcal C_{108}^{\rm red},
\qquad
R:\XTPK^{[108]}\longrightarrow\mathcal U_{12}^{\rm sgn}.
\]
L'application \(F\) ne retient que les champs du registre réduit, tandis que
\(R\) évalue la dodécaphase signée. L'existence d'une application \(\Xi\)
telle que \(R=\Xi\circ F\) équivaut à la factorisation de \(R\) par une
projection qui peut éliminer la coordonnée de l'observable, l'orientation,
la retenue ou la frontière. Cette factorisation constitue une proposition
supplémentaire et ne fait pas implicitement partie de la reconstruction sur l'espace
complet des états.

Le registre réduit est, par définition, une projection de l'état complet,
non l'état de départ. Les résultats ultérieurs utilisent \(R\) sur
l'état complet et ne présupposent pas de factorisation par \(F\).

\section{Systèmes inverses de survivants}

La transition TPK n'est pas nécessairement une fonction. Pour typer correctement
ce fait, soit \(V_t\) l'ensemble fini des états complets possibles à
la profondeur \(t\), et soit
\[
\mathcal R_t\subseteq V_t\times V_{t+1}
\]
la relation de transition admissible, avec tous les champs de l'état
enrichi conservés. Nous définissons alors \(X_t\), non comme l'ensemble des
états terminaux isolés, mais comme l'ensemble des préfixes admissibles, c'est-
à-dire des histoires admissibles complètes :
\[
X_t=
\Bigl\{
(v_0,\ldots,v_t)\in\prod_{j=0}^{t}V_j:
(v_j,v_{j+1})\in\mathcal R_j\ (0\le j<t),
\ \text{et les clôtures sont satisfaites jusqu’à }t
\Bigr\}.
\]
Sur ces objets, la troncature est bien une application canonique :
\[
\begin{aligned}
\rho_{t+1,t}:X_{t+1}&\longrightarrow X_t,\\
(v_0,\ldots,v_t,v_{t+1})&\longmapsto(v_0,\ldots,v_t).
\end{aligned}
\]
Elle est bien définie précisément parce que tout préfixe d'une histoire
admissible reste une histoire admissible pour les fermetures déjà imposées.
Si un filtre utilisait une information future et pouvait invalider
rétroactivement un préfixe, il faudrait d'abord remplacer \(X_t\) par le
sous-ensemble des préfixes prolongeables ; avec cette convention, la même troncature
redevient bien définie.

Une histoire HMT globale est une famille compatible de préfixes
\[
x=(x_t)_{t\ge0},
\qquad
x_t\in X_t,
\qquad
\rho_{t+1,t}(x_{t+1})=x_t.
\]
Elle équivaut à une suite infinie \((v_0,v_1,\ldots)\) dont les paires
consécutives appartiennent aux relations \(\mathcal R_t\) et dont les préfixes
satisfont toutes les fermetures. L'espace des histoires est la limite inverse
\[
X_\infty=\varprojlim X_t.
\]
Cette formulation évite d'imposer une transition locale univoque
\(B_t\mapsto B_{t+1}\). La transition peut être une relation finie, plusieurs
branches peuvent partager le même état terminal visible et l'unicité peut
n'apparaître qu'après avoir imposé la compatibilité future. En conservant le
préfixe complet, on ne perd pas prématurément l'information permettant
de les distinguer.

\begin{theorem}[Section coinductive unique]
Supposons que, pour tout \(t\), l'ensemble des survivants \(S_t\subseteq
X_t\) soit non vide, que les restrictions
\(\rho_{t+1,t}:S_{t+1}\to S_t\) soient compatibles et que chaque \(S_t\) contienne un
unique préfixe. Alors \(\varprojlim S_t\) contient une unique histoire
globale.
\end{theorem}

\begin{proof}
Soit \(x_t\) l'unique préfixe de \(S_t\). La compatibilité impose
\(\rho_{t+1,t}(x_{t+1})=x_t\), de sorte que \((x_t)_t\) appartient à la limite
inverse. Toute autre histoire aurait à chaque niveau un membre de \(S_t\)
et, par unicité, coïnciderait préfixe par préfixe avec \((x_t)_t\).
\end{proof}

Le théorème formalise le passage de la compatibilité niveau par niveau à une histoire
globale. Il ne transforme pas une fenêtre finie en une conclusion coinductive : son
application à une famille concrète exige de vérifier à chaque profondeur la non-
vacuité, la compatibilité des restrictions et l'unicité des préfixes
survivants.

\section{Connexion nonadique de survie}
\label{sec:conexion-nonadica}

La contraction des cylindres ne procède pas bloc par bloc comme une substitution
déterministe. À certaines frontières, il existe plusieurs répartitions locales d'une
même donnée agrégée ; la prolongeabilité décide laquelle appartient à l'histoire
publiée. La structure résultante est une connexion sur les neuf classes de
phase : à l'achèvement d'un tour, la classe de phase revient, mais l'état de la
fibre a changé. Cette section construit la fenêtre finie complète dans laquelle
ce retour est observé.

Le domaine de ce calcul local est le registre de transport engendré par la
composition APP--TRIT--TPK. Ses états initiaux, sa matrice entière de
propagation et ses relations de frontière sont des publications internes préalables de
la transition enrichie \(\mathcal U_t\), conformément au registre orienté de
douze transitions et à la reconstruction unique
\(L_i=X_i^{-1}Y_i\). Cette section les reçoit comme des résultats déjà démontrés,
non comme des prémisses externes ni comme des données conventionnelles. Tous ses énoncés
sont exacts sur ces sorties. Le théorème coinductif précédent étend la
même construction à une profondeur arbitraire et conserve la non-vacuité,
la compatibilité des troncatures et l'unicité des préfixes survivants.

\subsection{Préfixe de \(30\) trits et 1re frontière}

Nous écrivons \(u_k^\chi\in\FF_3^6\), avec
\(\chi\in\{\pi,e,\varphi\}\), pour le bloc visible du canal \(\chi\) à la
profondeur \(k\). Les cinq premiers blocs forment trois préfixes de trente
trits :
\begin{align*}
w_{30}^{\pi}&=010211|012222|010211|002111|110221,\\
w_{30}^{e}&=201101|121221|102011|012222|102011,\\
w_{30}^{\varphi}&=121200|112202|121020|010210|010200.
\end{align*}
Le sixième bloc est la première frontière postérieure à ces préfixes :
\[
R_{36}=
\begin{pmatrix}
u_5^\pi\\ u_5^e\\ u_5^\varphi
\end{pmatrix}
=
\begin{pmatrix}
222220\\021222\\102011
\end{pmatrix}.
\]
La notation \(R_{36}\) indique que trente-six trits par
canal ont été exposés. Elle ne désigne pas un état terminal : le registre reconstruit ci-dessous
contient vingt blocs, et la récurrence admet des continuations ultérieures
lorsque les états de Hensel correspondants sont fournis.

\subsection{Axe d'auto-échelle et fibre résiduelle}

\ifdefined\HMTAxialTheoremDefined
La décomposition axiale de la signature et l'identification de l'axe
\(\varphi\) face à la fibre résiduelle \(\pi/e\) ont déjà été démontrées dans le
\cref{thm:eje-espejo-nonadico}. On conserve ici leur conséquence pour la
fenêtre nonadique sans dupliquer le théorème ni la preuve.
\else
Pour
\[
B_k=(u_k^\pi,u_k^e,u_k^\varphi)\in(\FF_3^6)^3
\]
nous définissons, coordonnée par coordonnée,
\[
q_k=u_k^\pi+u_k^e+u_k^\varphi,
\qquad
a_k=u_k^\pi+u_k^e-u_k^\varphi.
\]

\begin{theorem}[Décomposition axiale de la signature de frontière]
\label{thm:eje-espejo-nonadico}
La paire \((q_k,a_k)\) détermine de manière univoque le canal axial et l'agrégat des
deux autres canaux :
\[
u_k^\varphi=2(q_k-a_k),
\qquad
u_k^\pi+u_k^e=q_k-u_k^\varphi.
\]
Par conséquent, une fibre de l'application
\[
q_{\rm ax}(u^\pi,u^e,u^\varphi)
=(u^\pi+u^e,u^\varphi)
\]
ne peut conserver d'ambiguïté que dans la répartition interne entre \(u^\pi\) et
\(u^e\).
\end{theorem}

\begin{proof}
La soustraction des deux équations de définition donne
\(q_k-a_k=2u_k^\varphi\). L'inverse de \(2\) dans \(\FF_3\) est encore
\(2\), d'où
\(u_k^\varphi=2(q_k-a_k)\). Soustraire ce vecteur de \(q_k\) produit
\(u_k^\pi+u_k^e\). L'opération ne retrouve pas chaque terme séparément ;
l'unique fibre résiduelle est donc l'ensemble des répartitions de l'agrégat fixé.
\end{proof}

Dans la terminologie géométrique du modèle, \(\varphi\) constitue l'axe
fixé par la signature et la paire \(\pi/e\) la fibre résiduelle. L'expression ne
postule pas de symétrie analytique des constantes réelles : elle nomme la
décomposition linéaire exacte qui vient d'être démontrée dans \(\FF_3^6\).
\fi

\subsection{Orbite complète des \(9\) phases consécutives du Topological Phase Kernel}

Nous fixons la phase absolue au moyen de
\[
g(k)=
\begin{cases}
k\bmod9,&k\not\equiv0\pmod9,\\
9,&k\equiv0\pmod9.
\end{cases}
\]
Comme \(w_{30}\) occupe \(k=0,\ldots,4\), le premier tour complet postérieur
au préfixe parcourt \(k=5,\ldots,13\), avec les phases
\(5,6,7,8,9,1,2,3,4\). L'énumération des \(729=3^6\) candidats de
signature agrégée à chaque niveau produit la fenêtre suivante :
\begin{center}
\small
\begin{tabular}{@{}ccrrlll@{}}
\toprule
\(k\)&\(g(k)\)&locaux&sélectionnés&\(u_k^\pi\)&\(u_k^e\)&\(u_k^\varphi\)\\
\midrule
5&5&1&1&222220&021222&102011\\
6&6&1&1&111201&201202&221021\\
7&7&3&1&212121&222210&200221\\
8&8&3&1&200121&212212&110110\\
9&9&2&1&100100&020112&010122\\
10&1&1&1&101222&112221&112002\\
11&2&1&1&022212&110001&100021\\
12&3&1&1&012012&202222&000120\\
13&4&1&1&111210&112102&211122\\
\bottomrule
\end{tabular}
\end{center}
La colonne « locaux » compte les états d'entrée admissibles au niveau
indiqué ; elle ne compte pas les sorties de la relation de frontière. Cette distinction
est essentielle dans \(k=9\) : les deux états locaux d'entrée précèdent une
fibre de trois sorties de l'état sélectionné.

\subsection{Retour 9→1 de la phase nonadique avec changement d'état et avancée de la mémoire}

Soit
\[
 \pi_B:V_k\longrightarrow\mathcal B:=(\FF_3^6)^3
\]
la projection sur le triplet visible. La relation de transport sur les états
complets induit, une fois fixées les données de fibre du registre, une relation
visible \(\MonNine^B\). L'état sélectionné qui entre dans la neuvième phase est
\[
B_9=(100100\mid020112\mid010122).
\]
Son image relationnelle se compose exactement de
\[
\begin{aligned}
B_{10}^{(0)}&=(101222\mid112221\mid112002),\\
B_{10}^{(1)}&=(102221\mid111222\mid112002),\\
B_{10}^{(2)}&=(111221\mid102222\mid112002).
\end{aligned}
\]
Les trois sorties ont
\[
q_{10}=022112,
\qquad a_{10}=101111,
\]
et, par le \cref{thm:eje-espejo-nonadico}, partagent
\[
u_{10}^\varphi=112002,
\qquad u_{10}^\pi+u_{10}^e=210110.
\]
Elles diffèrent exclusivement par les trois répartitions publiées de cet agrégat
entre \(\pi\) et \(e\).

\begin{theorem}[Sélection par deux niveaux de prolongeabilité]
\label{thm:seleccion-novena-fase}
Dans la fenêtre calibrée, les trois sorties précédentes admettent un prolongement
de profondeur un. En imposant également le bloc suivant
\[
B_{11}=(022212\mid110001\mid100021)
\]
et la frontière décimale \((502,662,638)\), seule
\(B_{10}^{(0)}\) admet un prolongement du même état complet. Par conséquent
\[
\#\operatorname{Surv}^{\rm cal}_{1}(B_9)=3,
\qquad
\#\operatorname{Surv}^{\rm cal}_{2}(B_9)=1.
\]
\end{theorem}

\begin{proof}
Les triplets décimaux associés aux trois sorties sont, respectivement,
\[
(279,352,365),\quad(280,351,365),\quad(282,350,365).
\]
Pour un mot ternaire de longueur \(L\), d'entier associé \(N\), et un
préfixe de \(b\) triplets, d'entier associé \(D\), la compatibilité des
deux cylindres équivaut aux deux inégalités entières
\[
N1000^b<(D+1)3^L,
\qquad
(N+1)1000^b>D3^L.
\]
Les six inégalités de chaque sortie sont satisfaites au niveau actuel. En
annexant \(B_{11}\) et \((502,662,638)\), elles sont à nouveau satisfaites pour les trois
canaux de \(B_{10}^{(0)}\). Pour \(B_{10}^{(1)}\) et \(B_{10}^{(2)}\),
l'inégalité du canal \(\varphi\) est conservée, mais celles de \(\pi\) et
\(e\) échouent. L'évaluation s'effectue avec des entiers dans le certificat fini correspondant ; aucune
tolérance numérique n'intervient. Par conséquent, l'ensemble des candidats
se réduit de trois à un lors du passage de la profondeur un à la profondeur deux.
\end{proof}

Les deux autres mots ne peuvent précéder une même queue ternaire que si
l'extension profonde de Hensel est modifiée. Ils ne sont donc pas des prolongements du
même état complet du registre : le reste visible coïncide en partie, mais
la généalogie de la fibre est distincte.

\subsection{Retour de phase avec changement d'état}

La phase satisfait \(g(k+9)=g(k)\). L'état visible, en revanche, ne se
réinitialise pas. En particulier,
\[
B_1=(012222\mid121221\mid112202),
\qquad
B_{10}=(101222\mid112221\mid112002),
\]
de sorte que \(g(1)=g(10)=1\) mais \(B_1\ne B_{10}\). La vérification des
neuf paires \(k\mapsto k+9\), pour \(k=1,\ldots,9\), donne en outre
\[
u_{k+9}^\varphi\ne u_k^\varphi,
\qquad
u_{k+9}^\pi+u_{k+9}^e\ne u_k^\pi+u_k^e.
\]
Par conséquent, la fenêtre satisfait exactement
\[
\boxed{g(k+9)=g(k),\qquad B_{k+9}\ne B_k
\quad(1\le k\le9).}
\]
Il s'agit d'un retour de phase avec holonomie visible non triviale. La dénomination
\emph{monodromie nonadique} désigne cette action de retour sur la
fibre enrichie ; elle n'introduit pas de dixième phase. Pour élever l'observation
finie à une monodromie globale, il faut construire l'action à toute profondeur
et démontrer la compatibilité coinductive exigée au début de la section.

\subsection{La branche de \(20\) blocs}

Le retour précédent ne clôt pas l'émission. Les vingt blocs que le
certificat reconstruit à partir des trois états entiers publiés sont
\begin{align*}
\pi:\;&010211|012222|010211|002111|110221|222220|111201|\\
&212121|200121|100100|101222|022212|012012|111210|\\
&121011|200220|120210|000101|022010|020111,\\[2pt]
e:\;&201101|121221|102011|012222|102011|021222|201202|\\
&222210|212212|020112|112221|110001|202222|112102|\\
&102010|022020|222002|012010|001112|022212,\\[2pt]
\varphi:\;&121200|112202|121020|010210|010200|102011|221021|\\
&200221|110110|010122|112002|100021|000120|211122|\\
&202200|202110|221000|100102|111122|020010.
\end{align*}
Ainsi, \(w_{30}\) est le préfixe et \(R_{36}\) la première frontière d'une branche
explicitement continuée ; aucun des deux objets ne remplace l'état
enrichi qui transporte cette continuation.

\subsection{Décodage exact de \(12\) triplets décimaux au moyen du lecteur nonadique}

Soit \(N_{\chi,k}\) l'entier représenté en base trois par les \(k\) premiers
blocs du canal \(\chi\). Avec \(k=13\), le mot a \(78\) trits et
définit
\[
D_{\chi,12}
=\left\lfloor\frac{1000^{12}N_{\chi,13}}{3^{78}}\right\rfloor.
\]
L'écriture de \(D_{\chi,12}\) en douze blocs de trois chiffres est
\[
\begin{array}{c|rrrrrrrrrrrr}
\pi&141&592&653&589&793&238&462&643&383&279&502&884\\
e&718&281&828&459&045&235&360&287&471&352&662&497\\
\varphi&618&033&988&749&894&848&204&586&834&365&638&117.
\end{array}
\]
L'identité est entière : elle équivaut à
\[
D_{\chi,12}3^{78}
\le 1000^{12}N_{\chi,13}
<(D_{\chi,12}+1)3^{78}.
\]
Par conséquent, les deux cylindres se coupent dans chaque canal et l'opération de
lecture produit exactement les trente-six chiffres publiés. Ce
résultat est une conséquence du registre engendré. Les constructions
historiques qui ont utilisé des développements conventionnels appartiennent exclusivement
à la comparaison et au contrôle ultérieur : elles n'interviennent ni dans la sélection
orbitale, ni dans la transition \(\mathcal U_t\), ni dans la reconstruction de
\(L_0,L_1\), ni dans le calendrier \(0011\). L'identité certifie la lecture
archimédienne des états produits par APP--TRIT--TPK et conserve intacte
leur provenance générative indépendante de valeurs cibles.

\subsection{Calendrier de réouverture des cylindres et périodicité du prolongement compatible}

La conversion d'un bloc de six trits en triplets décimaux a pour pente
moyenne
\[
\lambda=6\log_{1000}3\in(0,1),
\qquad k(t)=\lfloor\lambda(t+1)\rfloor.
\]
Un indice \(t\) est une réouverture sans nouveau triplet lorsque
\(k(t+1)=k(t)\). Soit \(\delta=1-\lambda\).

\begin{theorem}[Calendrier de Beatty de la lecture ternaire--décimale]
\label{thm:calendario-beatty}
L'ensemble des réouvertures est
\[
\mathcal T=
\left\{\left\lceil\frac{m}{\delta}\right\rceil-2:m\ge1\right\}.
\]
Ses premiers éléments sont
\[
20,42,64,86,108,130,151,173,195,217,\ldots,
\]
et deux réouvertures consécutives sont distantes de \(21\) ou \(22\) blocs.
\end{theorem}

\begin{proof}
La quantité \(\lambda=\log_{1000}729\) est irrationnelle : une égalité
\(\lambda=p/q\) impliquerait \(729^q=1000^p\), ce qui est impossible par factorisation
en nombres premiers. \(\delta\) est également irrationnelle, d'où
\[
k(t)=(t+1)-\lceil\delta(t+1)\rceil.
\]
L'égalité \(k(t+1)=k(t)\) se produit exactement lorsque la fonction plafond du
second terme augmente d'une unité. La \(m\)-ième occurrence s'obtient
en résolvant
\(\delta(t+1)<m<\delta(t+2)\), ce qui donne
\(t=\lceil m/\delta\rceil-2\). Les différences d'une suite de Beatty
appartiennent à
\[
\left\{\left\lfloor\delta^{-1}\right\rfloor,
\left\lceil\delta^{-1}\right\rceil\right\}=\{21,22\}.\qedhere
\]
\end{proof}

\subsection{Récurrence de Bockstein--Hensel du reste et du quotient}

La mémoire transportée par la branche provient de la suite exacte
\[
0\longrightarrow\ZZ^6\xrightarrow{\,\times3\,}\ZZ^6
\longrightarrow\FF_3^6\longrightarrow0
\]
et de la matrice unimodulaire
\[
A=
\begin{pmatrix}
2&2&2&1&2&1\\
2&1&2&2&1&1\\
1&0&0&1&0&2\\
0&0&1&1&1&0\\
2&2&2&0&2&1\\
2&1&2&1&0&0
\end{pmatrix},
\qquad\det A=1.
\]
Pour un état ligne \(x_k^\chi\in\ZZ^6\), on définit
\[
y_k^\chi=x_k^\chi A,
\qquad
u_k^\chi=y_k^\chi\bmod3\in\{0,1,2\}^6,
\qquad
x_{k+1}^\chi=\frac{y_k^\chi-u_k^\chi}{3}.
\]

\begin{proposition}[Conservation exacte de la mémoire du quotient]
\label{prop:memoria-bockstein-hensel}
\(x_k^\chi\) et \(A\) étant fixés, la paire
\((u_k^\chi,x_{k+1}^\chi)\) est unique et satisfait
\[
x_k^\chi A=u_k^\chi+3x_{k+1}^\chi.
\]
Réciproquement, comme \(A\in\operatorname{GL}_6(\ZZ)\), la paire détermine
de manière univoque \(x_k^\chi\).
\end{proposition}

\begin{proof}
La division euclidienne de chaque coordonnée de \(x_k^\chi A\) par trois produit
un reste unique dans \(\{0,1,2\}\) et un quotient entier unique. Cela prouve
le premier énoncé et l'identité. Puisque \(\det A=1\), l'inverse
\(A^{-1}\) a des entrées entières ; par conséquent
\[
x_k^\chi=(u_k^\chi+3x_{k+1}^\chi)A^{-1}
\]
retrouve l'état précédent sans ambiguïté.
\end{proof}

Le reste \(u_k^\chi\) est la partie visible et le quotient
\(x_{k+1}^\chi\) la mémoire 3-adique transportée. Le certificat reproduit
soixante pas coordonnés ---vingt par canal---, vérifie \(AA^{-1}=I\) par
arithmétique entière et confronte tous les blocs au registre. La dénomination
Bockstein--Hensel fait ici référence à cette séparation reste--quotient et à son
itération 3-adique ; elle ne s'identifie pas sans données supplémentaires à un opérateur de
Bockstein cohomologique.

\section{Du préfixe à un point réel}

Associons à chaque préfixe compatible \(x_t\in X_t\) un intervalle fermé
\(I_t(x_t)\subset\mathbb R\). L'interprétation géométrique est celle d'un
cylindre : tous les futurs qui partagent le même préfixe se projettent à l'intérieur
de \(I_t\).

\begin{theorem}[Convergence archimédienne d'intervalles emboîtés]
Soit \((x_t)_{t\ge0}\) une histoire compatible. Si
\[
I_{t+1}(x_{t+1})\subseteq I_t(x_t)
\]
et
\[
\operatorname{diam} I_t(x_t)\longrightarrow0,
\]
alors il existe un unique \(r\in\mathbb R\) tel que
\[
\bigcap_{t\ge0}I_t(x_t)=\{r\}.
\]
\end{theorem}

\begin{proof}
C'est le principe des intervalles emboîtés. Les extrémités gauches forment une
suite croissante bornée et les droites une suite décroissante bornée. Leurs limites
existent ; leur différence est bornée par le diamètre, qui tend vers
zéro. Les deux limites coïncident en un point unique contenu dans tous les
intervalles.
\end{proof}

Soit \(X_\infty^{\rm Arch}\subseteq X_\infty\) le domaine des histoires pour
lesquelles les cylindres sont emboîtés et leurs diamètres tendent vers zéro. Si ces
deux propriétés font partie de l'admissibilité globale, alors
\(X_\infty^{\rm Arch}=X_\infty\) ; sinon, cette restriction de domaine est
nécessaire. Nous définissons l'évaluation
\[
\operatorname{ev}_{\mathbb R}:X_\infty^{\rm Arch}\longrightarrow\mathbb R,
\qquad
\operatorname{ev}_{\mathbb R}(x)=
\text{l’unique point de }\bigcap_{t\ge0}I_t(x_t).
\]
L'unicité de la valeur n'impose pas l'unicité de l'histoire. Deux généalogies
\(x\ne y\) peuvent satisfaire
\[
\operatorname{ev}_{\mathbb R}(x)
=\operatorname{ev}_{\mathbb R}(y).
\]
Nous définissons l'équivalence de lecture
\[
x\sim_{\rm ev}y
\quad\Longleftrightarrow\quad
\operatorname{ev}_{\mathbb R}(x)
=\operatorname{ev}_{\mathbb R}(y).
\]

\begin{theorem}[Quotient de l'évaluation archimédienne]
L'application
\[
\begin{aligned}
\overline{\operatorname{ev}}_{\mathbb R}:
X_\infty^{\rm Arch}/\!\sim_{\rm ev}
&\longrightarrow
\operatorname{Im}(\operatorname{ev}_{\mathbb R})\subseteq\mathbb R,\\
[x]&\longmapsto\operatorname{ev}_{\mathbb R}(x)
\end{aligned}
\]
est une bijection. En particulier, le quotient s'identifie à tout
\(\mathbb R\) si et seulement si l'évaluation
\(\operatorname{ev}_{\mathbb R}\) est surjective.
\end{theorem}

\begin{proof}
La définition de \(\sim_{\rm ev}\) rend la valeur indépendante du
représentant, de sorte que l'application est bien définie. Elle est injective car
deux classes ayant la même valeur sont, par définition, une seule classe ; elle est
surjective sur l'image par la définition même de l'image. Le dernier
énoncé équivaut à
\(\operatorname{Im}(\operatorname{ev}_{\mathbb R})=\mathbb R\).
\end{proof}

Le théorème de couverture du lecteur archimédien construit précisément ces
préfixes : pour tout intervalle compact \(K\subset\mathbb R\),
\(\operatorname{ev}_{\mathbb R}(\Omega_{\infty,K})=K\), et les intervalles
\([-m,m]\) forment une famille dirigée dont l'union est \(\mathbb R\). Par
conséquent, chaque réel possède au moins une histoire HMT compatible et
l'énoncé démontré est
\[
X_\infty^{\rm Arch}/\!\sim_{\rm ev}
\cong\operatorname{Im}(\operatorname{ev}_{\mathbb R})
=\mathbb R.
\]
La valeur archimédienne est donc le quotient qui identifie les fibres de
l'évaluation ; le nombre HMT enrichi conserve en outre la classe de
trajectoire, la retenue et l'orientation. L'identification ne découle pas de la
simple formation abstraite du quotient : elle découle du théorème de couverture déjà
démontré dans la construction discrète du continu.

La descente algébrique est un autre problème et doit être résolue opération par
opération. Supposons que deux opérations candidates aient été construites,
\(\oplus\) et \(\odot\), sur \(X_\infty^{\rm Arch}\), internes à ce
domaine. La somme passe au quotient si et seulement si
\[
x\sim_{\rm ev}x',\ y\sim_{\rm ev}y'
\Longrightarrow
x\oplus y\sim_{\rm ev}x'\oplus y',
\]
tandis que le produit passe au quotient, de manière indépendante, si et seulement si
\[
x\sim_{\rm ev}x',\ y\sim_{\rm ev}y'
\Longrightarrow
x\odot y\sim_{\rm ev}x'\odot y'.
\]
La première congruence n'implique pas la seconde. En outre, pour que les opérations
induites soient précisément la somme et le produit ordinaires du quotient
archimédien, il faut prouver séparément
\[
\operatorname{ev}_{\mathbb R}(x\oplus y)
=\operatorname{ev}_{\mathbb R}(x)+\operatorname{ev}_{\mathbb R}(y),
\]
\[
\operatorname{ev}_{\mathbb R}(x\odot y)
=\operatorname{ev}_{\mathbb R}(x)\operatorname{ev}_{\mathbb R}(y),
\]
et que l'image de l'évaluation soit stable par l'opération correspondante.
Ce n'est qu'après ces certificats que le quotient hérite de la structure de
sous-anneau ---ou de sous-corps, si les opposés et les inverses sont également construits---
de \(\mathbb R\). La bijection ensembliste du théorème précédent est formelle ;
le contenu spécifique de HMT consiste à construire l'évaluation, à caractériser ses
fibres, à démontrer la couverture et à établir chaque loi de descente.

\section{Représentation de Weyl--Heisenberg associée au tore nonadique}

La forme d'aire orientée du tore nonadique est
\[
\sigma((a,b),(c,d))=ad-bc\pmod9.
\]
Cette forme alternée est celle que l'on obtient, avec l'orientation choisie, en
intégrant la courbure mixte d'APP sur des parallélogrammes du tore. Soit
\(\omega=\exp(2\pi i/9)\) et
\(\mathcal H_9=\mathbb C^9\), de base \(\{|n\rangle:n\in\mathbb Z_9\}\).
Nous définissons
\[
X|n\rangle=|n+1\rangle,
\qquad
Z|n\rangle=\omega^n|n\rangle.
\]
Alors
\[
ZX=\omega XZ,
\]
et, plus généralement,
\[
(X^aZ^b)(X^cZ^d)
=\omega^{bc}X^{a+c}Z^{b+d}.
\]
Par conséquent
\[
 c\bigl((a,b),(c,d)\bigr)=bc\pmod9
\]
est le \(2\)-cocycle polarisé qui apparaît dans cette section de l'extension
centrale. Nous définissons explicitement
\[
 \operatorname{Heis}_9
 =\{\omega^kX^aZ^b:a,b,k\in\ZZ/9\ZZ\}.
\]
C'est un groupe d'ordre \(9^3=729\). Il s'agit d'une extension centrale associée
au groupe des translations du support \(X_9\) d'APP, non d'un sous-groupe du
graphe ni de la catégorie des chemins d'APP\@.
Le commutateur des deux transports est
\[
(X^aZ^b)(X^cZ^d)(X^aZ^b)^{-1}(X^cZ^d)^{-1}
=\omega^{bc-ad}I.
\]
Par conséquent, son exposant est \(-\sigma((a,b),(c,d))\). Le cocycle
polarisé \(c\) et la forme \(\sigma\) ne sont pas le même objet : la seconde est,
avec cette convention, l'opposée de l'antisymétrisation
\(c(u,v)-c(v,u)\). La phase du commutateur est ainsi l'inverse de la phase de
Wilson pour l'orientation fixée.

Pour construire des observables, nous prenons
\[
Q_{\rm H}=\operatorname{diag}(0,1,\ldots,8),
\qquad
F_{jk}=\frac13\omega^{jk},
\qquad
P_{\rm H}=F^\dagger Q_{\rm H}F.
\]
Les opérateurs \(Q_{\rm H},P_{\rm H}\) sont autoadjoints. Leurs bases propres sont mutuellement
non biaisées car \(|F_{jk}|^2=1/9\). Pour tout vecteur unitaire
\(\psi\in\mathcal H_9\), l'inégalité de Robertson donne
\[
\operatorname{Var}_\psi(Q_{\rm H})\,
\operatorname{Var}_\psi(P_{\rm H})
\ge
\frac14\left|\langle\psi,[Q_{\rm H},P_{\rm H}]\psi\rangle\right|^2,
\]
et la borne entropique pour les bases mutuellement non biaisées produit
\[
H_{Q_{\rm H}}(\psi)+H_{P_{\rm H}}(\psi)\ge\log9.
\]
Ici \(H\) est l'entropie de Shannon, avec logarithme naturel, des
distributions spectrales correspondantes. Les relations construites se
résument, sans leur attribuer de causalité physique supplémentaire, par
\[
\begin{gathered}
\text{forme alternée }(X_9,\sigma)
\longrightarrow
\operatorname{Heis}_9,\\
\operatorname{Heis}_9\longrightarrow(X,Z)
\longrightarrow(Q_{\rm H},P_{\rm H})
\longrightarrow\text{incertitude finie}.
\end{gathered}
\]
Il s'agit d'un modèle algébrique fini d'opérateurs ; il ne spécifie pas à lui seul
une dynamique physique ni des unités de mesure.

\section{Dynamique cyclique sur \texorpdfstring{\(\mathbb Z_9\times\mathbb Z_{12}\)}{Z9 x Z12} et quotient prédictif minimal}

Le système dynamique fini \(\mathfrak C_{9,12}\) permet de déterminer la mémoire
minimale nécessaire pour que l'évolution observée soit déterministe. Soit
\[
X=\{(b,m):b\in\mathbb Z/9\mathbb Z, 0\le m<12\}
\]
avec successeur
\[
S(b,m)=
\begin{cases}
(b,m+1),&m<11,\\
(b+1\bmod9,0),&m=11.
\end{cases}
\]
Les étiquettes marginales sont
\[
\beta(b,m)=4b\pmod{12},
\qquad
d(b,m)=1+(m+\beta(b,m)\pmod{12}).
\]
Pour une application d'observation \(q\), le mot futur inclusif d'horizon \(h\) est
\[
Q_h^q(x)=\bigl(q(x),q(Sx),\ldots,q(S^hx)\bigr).
\]

\begin{theorem}[Quotient prédictif minimal]
Soient \(X\) fini, \(T:X\to X\) et \(q:X\to V\). Définissons
\[
\operatorname{Eq}(f):=\{(x,y)\in X^2:f(x)=f(y)\},
\qquad
E_h=\operatorname{Eq}(Q_h^q),
\qquad
E_\infty=\bigcap_{k\ge0}\operatorname{Eq}(q\circ T^k).
\]
Alors
\(E_{h+1}=E_h\cap(T\times T)^{-1}(E_h)\), la suite se stabilise en temps
fini et \(X/E_\infty\) est le quotient ayant le plus petit nombre d'états qui
raffine \(q\) et admet une dynamique déterministe induite.
\end{theorem}

\begin{proof}
\(E_h\) exige l'égalité des observations aux temps \(0,\ldots,h\), tandis que
\((T\times T)^{-1}(E_h)\) exige l'égalité en \(1,\ldots,h+1\). Leur intersection est
\(E_{h+1}\). Comme \(X\) est fini, le nombre de classes ne peut pas croître
indéfiniment. Au premier niveau stable, \(T\) respecte les classes et
passe au quotient. Toute équivalence \(T\)-compatible contenue dans
\(\operatorname{Eq}(q)\)
n'identifie que des états ayant le même futur complet, de sorte qu'elle est
contenue dans \(E_\infty\) ; d'où la propriété minimale.
\end{proof}

L'évaluation exhaustive de \(\mathfrak C_{9,12}\) donne :
\begin{center}
\small
\begin{tabular}{@{}lrrrr@{}}
\toprule
observable & valeurs instantanées & indice \(h\) & remarques & quotient stable\\
\midrule
\(\beta\) & 3 & 11 & 12 & \(C_{36}\), fibra 3\\
\(d\) & 12 & 8 & 9 & \(C_{36}\), fibra 3\\
\((\beta,d)\) & 36 & 0 & 1 & \(C_{36}\), fibra 3\\
\bottomrule
\end{tabular}
\end{center}
Dans les deux observables marginaux,
\[
\operatorname{Eq}(Q_{11}^{\beta})
=\operatorname{Eq}(Q_8^d)
=\operatorname{Eq}(\pi_{36}),
\qquad
\pi_{36}(b,m)=(b\bmod3,m).
\]
Le facteur \(36\) n'est pas introduit comme objectif. Il apparaît comme le quotient
prédictif minimal de chaque marginal. La phase \(m\) est une coordonnée d'état
nécessaire pour prédire ces futurs ; cet énoncé mathématique ne présuppose pas
d'interprétation causale physique. La transformation \(b\mapsto b+3\) est une
symétrie invisible pour ces observables à tout instant.

En tant que système dynamique abstrait, le quotient est un cycle de trente-six
états. Écrire \(C_{36}\) ne choisit toutefois pas d'origine distinguée :
pour obtenir une coordonnée \(\mathbb Z/36\mathbb Z\), il faut la fixer. Cela ne
transforme pas non plus automatiquement ces états en les trente-six paires d'une
nonade. En effet, \(\operatorname{Aut}J(9,2)\cong S_9\). Si les deux points d'une
paire sont dans des cycles distincts de longueurs \(r,s\) d'une permutation,
l'orbite de la paire est de longueur \(\operatorname{mcm}(r,s)\), avec \(r+s\leq9\),
et donc de longueur au plus \(20\). S'ils sont dans un même cycle de longueur
\(r\leq9\), l'orbite est de longueur \(r\), sauf dans le cas antipodal, où elle
est de longueur \(r/2\). Aucun automorphisme n'induit donc d'orbite de longueur
\(36\) sur les paires. Par conséquent, le passage
\[
C_{36}\longleftrightarrow \binom{[9]}2
\]
reste une identification calibrée, non une identification équivariante naturelle.
Même un ordre hamiltonien des paires prouverait une conservation locale
de l'incidence pas à pas, non un automorphisme global ni la canonicité du
dictionnaire.

\section{Nombre HMT comme histoire compatible dans \(X_\infty\) : identité généalogique, structure de lecture et évaluation archimédienne}

Les résultats précédents déterminent l'instance suivante de la définition
directrice. Dans la notation de cette section, un nombre HMT est une histoire
compatible
\[
x\in X_\infty=\varprojlim X_t.
\]
Un lecteur \(L\) est une structure supplémentaire : la paire \((x,L)\) constitue une
présentation lue ou une mesure, mais ne modifie pas le type d'identité de
\(x\). Lorsque \(x\) appartient à \(X_\infty^{\rm Arch}\), sa valeur
archimédienne est l'intersection ponctuelle des cylindres ; l'ensemble des
valeurs ainsi construit est
\(\operatorname{Im}(\operatorname{ev}_{\mathbb R})\subseteq\mathbb R\), et
son identification à toute la droite exige le théorème supplémentaire de couverture.
Sa mémoire minimale relative à un observable est le quotient prédictif des
futurs ; et la forme alternée du support somme--produit admet la
représentation finie de Weyl--Heisenberg construite ci-dessus. La valeur, la classe de trajectoire et l'incertitude sont
définies par trois opérations sur des domaines spécifiés d'une même
architecture de transport.
